\chapter{Procedure}

\section{Panel efficiency}
The efficiency of the panels can be calculated in the following way: All events in which e.g. scintillator 1 and 3 were triggered would suggest that scintillator 2 should have been triggered as well. This works fine for scintillators 2 and 3. For scintillators 1 and 4 the method has to be adjusted a bit. For e.g. scintillator 1 all events are taken into account in which at least two of the other scintillators were triggered. The problem is that this could have been triggered by a muon which is approaching at a slight angle in such a way that in never passed through scintillator 1. For scintillator 4 the decay of the muon is another effect which slightly changes the efficiency. These effects are not taken into account in this experiment. 

First the data is inspected to make sure all panels and the trigger work properly.
For the scintillators the following efficiencies where calculated
\begin{table}[h]
    \centering
    \begin{tabular}{c|c}
         scintillator& efficiency \\ 
         1 & 0.994 \\
         2 & 0.941 \\
         3 & 0.952 \\
         4 & 0.945 \\
    \end{tabular}
    \caption{Caption}
    \label{tab:my_label}
\end{table}
